\begin{abstract}
Catastrophic forgetting (CF) remains a central obstacle to deploying
Large Language Models (LLMs) in settings where new tasks arrive
sequentially. While recent work has produced a broad zoo of
mitigation strategies---ranging from elastic-weight regularization
and experience replay to orthogonal low-rank adapters---almost all
evaluations are performed on 7B--13B-parameter models trained on
data-center-grade hardware. In this work we ask a different but
equally important question: \emph{which mitigation strategies remain
effective when the practitioner is restricted to a sub-2B parameter
model and a single free-tier GPU?} We present a systematic empirical
comparison of six representative strategies---full fine-tuning,
vanilla LoRA, LoRA with Elastic Weight Consolidation (EWC),
LoRA with Experience Replay, Orthogonal LoRA (O-LoRA), and a
\emph{Minimal Mixed Rehearsal} (MMR) recipe we formalize from a
recent empirical observation---applied to Qwen2.5-1.5B-Instruct and
Phi-3.5-mini-instruct across a four-task sequence drawn from GLUE
(SST-2 $\rightarrow$ MRPC $\rightarrow$ CoLA $\rightarrow$ MNLI).
All experiments are run under QLoRA 4-bit quantization on a single
NVIDIA T4. We report Average Accuracy, Forgetting, Backward Transfer,
wall-clock training time, and peak GPU memory, and visualize the
resulting accuracy/compute trade-off as a Pareto frontier.
Our results show that \todo{insert one-sentence headline finding once
results are in}, and that MMR with a replay fraction as low as
\todo{X\%} matches or exceeds heavier methods while incurring
negligible additional compute. We release all code, configurations,
and run logs to enable reproduction on free-tier hardware.
\end{abstract}
